\section{The Raghavan--Xiao basis and the countable-basis problem}
\label{sec:rxcb}

Throughout this section, if $D$ and $H$ are Borel digraphs, then
$D\preceq_{\mathrm B}H$ means that there is a Borel homomorphism from
$D$ to $H$.  We write $D\preceq^{\min}_{\mathrm B}H$ when the
homomorphism can be chosen to be minimal in the sense of
\cite[Definition~3.10]{RX24}.  Let $\mathcal C_\infty$ denote the class
of Borel digraphs of uncountable Borel dichromatic number.  The
citation key \texttt{RX24} is a placeholder for the Raghavan--Xiao
entry in the surrounding bibliography.

\subsection{What is, and is not, proved in \cite{RX24}}
\label{subsec:rxcb-audit}

Fix a dense selector $E$.  For $f\in\omega^\omega$ satisfying
$f(n)\geq2$ for every $n$, put
\[
  \mathcal G_f=\mathcal G_0(E,f).
\]
Theorem~3.23 of \cite{RX24} says that, for every
$D\in\mathcal C_\infty$, there are $f$ and a continuous minimal
homomorphism $\mathcal G_f\to D$.  Theorem~3.25 gives the corresponding
statement for analytic targets with ``continuous homomorphism'' in
place of ``continuous minimal homomorphism.''  Conversely,
Lemma~3.8 implies that every $\mathcal G_f$ belongs to
$\mathcal C_\infty$.  Consequently,
\begin{equation}
  D\in\mathcal C_\infty
  \quad\Longleftrightarrow\quad
  (\exists f)\bigl(\mathcal G_f\preceq_{\mathrm B}D\bigr).
  \label{eq:rxcb-basis}
\end{equation}

The paper proves the upper bound that a uniformly real-parametrized
family, and hence a family of cardinality $\mathfrak c$, is a basis.
It does not prove a lower bound on the cardinality of an arbitrary
basis.  In particular, it does not rule out a finite or countable
basis under ordinary Borel homomorphism.  The displayed question
immediately following the digraph dichotomies is Question~3.26, which
asks whether minimality can also be obtained for analytic non-Borel
targets.  Questions~3.42 and~4.23 concern, respectively, the
classification of quasi-orders of countable Borel order dimension and
a possible ZFC upper bound for the Borel order dimension of locally
countable quasi-orders.  None asks for, or proves, the least size of a
basis for $\mathcal C_\infty$.

The paper also does not state that the fixed-$E$ family
$\{\mathcal G_f:f\in(\omega\setminus2)^\omega\}$, or any cofinal
subfamily of it, is pairwise incomparable under ordinary Borel
homomorphism.  This point is logically important; see
Lemma~\ref{lem:rxcb-antichain-basis} below.

\subsection{The descriptive complexity of the canonical basis}
\label{subsec:rxcb-complexity}

\begin{proposition}
\label{prop:rxcb-basis-complexity}
Suppose that $E$ is a recursive dense selector.  In any standard
effective coding of Borel digraphs, the map
\[
  f\longmapsto \operatorname{code}(\mathcal G_0(E,f))
\]
can be chosen $\Delta^1_1$.  Its untagged range is therefore
$\boldsymbol\Sigma^1_1$.  If $f$ is retained as part of the code, the
resulting tagged family is Borel, and it can be made $\Delta^1_1$ in a
suitable effective coding.  In particular, the Raghavan--Xiao basis
has an analytic, indeed naturally Borel, presentation and is not
$\boldsymbol\Pi^1_2$-hard.
\end{proposition}

\begin{proof}
Put
\[
  X_f=\prod_{n\in\omega}f(n).
\]
Uniformly in $f$, this is a closed subset of $\omega^\omega$.  For
$x,y\in\omega^\omega$, one has
\begin{align*}
 (x,y)\in R_{E,f}
 \quad\Longleftrightarrow\quad
 &x,y\in X_f\ \text{ and there is }\ell\in\omega\text{ such that}\\
 &x\mathbin{\upharpoonright}\ell
   =y\mathbin{\upharpoonright}\ell
   =E(f\mathbin{\upharpoonright}\ell),\\
 &y(\ell)=x(\ell)+1\pmod{f(\ell)},\quad\text{and}\quad
   (\forall k>\ell)\,x(k)=y(k).
\end{align*}
For fixed $\ell$, this is a closed condition on $(f,x,y)$.  The
quantification over $\ell$ is a countable union, so the universal arc
relation is uniformly $\boldsymbol\Sigma^0_2$.  A Borel code for the
closed vertex set and the $F_\sigma$ arc relation is therefore
obtained effectively from $f$.  This gives the asserted
$\Delta^1_1$ coding map.  Its range is analytic.  Retaining the
parameter $f$ gives an injective Borel parametrization, whose range is
Borel by the Lusin--Souslin theorem.
\end{proof}

\begin{corollary}
\label{cor:rxcb-high-delta12}
Relative to a standard space of valid Borel-digraph codes,
$\mathcal C_\infty$ is both $\boldsymbol\Sigma^1_2$ and
$\boldsymbol\Pi^1_2$, and hence is $\boldsymbol\Delta^1_2$.
\end{corollary}

\begin{proof}
By \eqref{eq:rxcb-basis}, a code $d$ belongs to
$\mathcal C_\infty$ if and only if there are a real $f$ and a code $u$
for a Borel function such that $u$ is a homomorphism
$\mathcal G_f\to D_d$.  Once $f,u,d$ are fixed, totality and
preservation of arcs form a $\boldsymbol\Pi^1_1$ condition on the
codes.  Pairing the existentially quantified reals gives a
$\boldsymbol\Sigma^1_2$ definition.

On the other hand, countable Borel dichromatic number is witnessed by
one real coding a sequence of Borel acyclic sets covering the vertex
space.  Validity of the cover and acyclicity of every member are
$\boldsymbol\Pi^1_1$ conditions.  Thus the complement of
$\mathcal C_\infty$ is $\boldsymbol\Sigma^1_2$, so
$\mathcal C_\infty$ is $\boldsymbol\Pi^1_2$ as well.
\end{proof}

The complexity proof used at a finite dichromatic threshold therefore
does not transfer to the uncountable threshold.  A countable basis
would yield a $\boldsymbol\Sigma^1_2$ definition, but the
Raghavan--Xiao basis already yields such a definition.  This does not
rule out every conceivable use of descriptive complexity, but it
removes the $\boldsymbol\Sigma^1_2$ versus
$\boldsymbol\Pi^1_2$-hard separation that drives the finite-threshold
argument.

\subsection{The consequence of a genuine antichain basis}
\label{subsec:rxcb-antichain}

\begin{lemma}[Antichain-basis lemma]
\label{lem:rxcb-antichain-basis}
Let $(\mathcal C,\preceq)$ be a quasi-ordered class.  Suppose that
$\mathcal A\subseteq\mathcal C$ is both a basis for $\mathcal C$ and
an antichain under $\preceq$.  Then every basis $\mathcal B$ for
$\mathcal C$ satisfies
\[
  |\mathcal B|\geq|\mathcal A|.
\]
\end{lemma}

\begin{proof}
For each $A\in\mathcal A$, choose $B_A\in\mathcal B$ with
$B_A\preceq A$.  Since $B_A\in\mathcal C$ and $\mathcal A$ is a
basis, choose $A'\in\mathcal A$ with $A'\preceq B_A$.  Then
$A'\preceq A$, so incomparability gives $A'=A$.  Hence
\begin{equation}
  A\preceq B_A\preceq A.
  \label{eq:rxcb-sandwich}
\end{equation}
If $A_0\neq A_1$ and $B_{A_0}=B_{A_1}$, then
\eqref{eq:rxcb-sandwich} gives both $A_0\preceq A_1$ and
$A_1\preceq A_0$, contradicting that $\mathcal A$ is an antichain.
Thus $A\mapsto B_A$ is injective.
\end{proof}

\begin{corollary}
\label{cor:rxcb-conditional-full}
If a continuum-sized Raghavan--Xiao family is simultaneously
\begin{enumerate}
  \item a basis for $\mathcal C_\infty$ under ordinary Borel
  homomorphism, and
  \item pairwise incomparable under that same notion of homomorphism,
\end{enumerate}
then every basis for $\mathcal C_\infty$ has cardinality
$\mathfrak c$.  In particular, there is no countable basis.
\end{corollary}

\begin{proof}
Apply Lemma~\ref{lem:rxcb-antichain-basis} with
$\preceq=\preceq_{\mathrm B}$.
\end{proof}

Thus, if the private statement that the members of ``the basis'' are
pairwise non-homomorphic refers to an actual cofinal family and to
ordinary Borel homomorphisms, the conjecture is already an immediate
corollary.  That interpretation is incompatible with the simultaneous
assertion that the least size of a basis is open.  The private
statement may instead concern a non-cofinal subfamily, a different
choice of parameters, or minimal homomorphisms.  The published paper
itself does not supply the incomparability hypothesis needed in
Corollary~\ref{cor:rxcb-conditional-full}.

\subsection{Closed walks and the first diagonal obstruction}
\label{subsec:rxcb-walks}

For an edge $(x,y)$ of $\mathcal G_g$, its \emph{level} is the unique
$m$ at which $x$ and $y$ differ.  For
$B\subseteq\omega\setminus2$, let
\[
  \langle B\rangle_+
  =\left\{\sum_{i<r}q_i b_i:
    r\geq1,\ b_i\in B,\ q_i\in\omega\setminus\{0\}\right\}
\]
be the positive additive monoid generated by $B$.

\begin{lemma}[Level counting]
\label{lem:rxcb-level-counting}
Let $W$ be a directed closed walk of $k$ edges in $\mathcal G_g$, and
let $N_m(W)$ be the number of edges of $W$ having level $m$.  Then
\[
  g(m)\mid N_m(W)
\]
for every $m$.  Consequently,
\[
  k\in\langle\operatorname{ran}(g)\rangle_+,
\]
and every level $m$ used by $W$ satisfies $g(m)\leq k$.
\end{lemma}

\begin{proof}
An edge of level $m$ increases coordinate $m$ by one modulo $g(m)$.
Every edge of a different level leaves coordinate $m$ unchanged.
Since the walk returns to its initial vertex, its total increment in
coordinate $m$ is zero modulo $g(m)$.  Hence $g(m)$ divides $N_m(W)$.
Summing the nonzero values $N_m(W)$ gives the other conclusions.
\end{proof}

\begin{corollary}
\label{cor:rxcb-monoid-obstruction}
If there is a homomorphism, with no definability assumption, from
$\mathcal G_f$ to $\mathcal G_g$, then
\[
  \operatorname{ran}(f)
  \subseteq\langle\operatorname{ran}(g)\rangle_+.
\]
Consequently, given $\langle f_n:n\in\omega\rangle$, if there is an
infinite $B\subseteq\omega\setminus2$ such that
\begin{equation}
  \operatorname{ran}(f_n)\not\subseteq\langle B\rangle_+
  \quad\text{for every }n,
  \label{eq:rxcb-monoid-diagonal}
\end{equation}
then every one-to-one enumeration $g$ of $B$ satisfies
$\mathcal G_{f_n}\not\preceq_{\mathrm B}\mathcal G_g$ for every $n$.
\end{corollary}

\begin{proof}
Every edge at level $\ell$ of $\mathcal G_f$ belongs to a canonical
cycle with exactly $f(\ell)$ vertices and edges.  A homomorphism sends
this cycle to a directed closed walk of $f(\ell)$ edges.  Apply
Lemma~\ref{lem:rxcb-level-counting}.
\end{proof}

Condition \eqref{eq:rxcb-monoid-diagonal} cannot always be obtained.
For $k\geq2$, let
\[
  f_k(n)=k(n+1).
\]
If $B$ is nonempty and $b\in B$, then
\[
  \operatorname{ran}(f_b)=\{b,2b,3b,\dots\}
  \subseteq\langle B\rangle_+.
\]
Thus a diagonalization which remembers only the lengths of closed
walks cannot settle the ordinary-homomorphism problem.  Ordinary
homomorphisms may fold one cycle onto a shorter cycle, or map it to a
concatenation of target cycles; for example, the natural map
$C_6\to C_3$ is a homomorphism.

There is a stronger tail version for continuous homomorphisms.

\begin{lemma}[Tail-monoid obstruction]
\label{lem:rxcb-tail-monoid}
If there is a continuous homomorphism
$\Phi:\mathcal G_f\to\mathcal G_g$, then, for every $N$, there is an
$L$ such that
\[
  f(\ell)\in
  \left\langle\{g(m):m\geq N\}\right\rangle_+
  \quad\text{for every }\ell\geq L.
\]
\end{lemma}

\begin{proof}
The space $X_f$ is compact.  By uniform continuity, for each $N$ there
is $L$ such that
\[
  x\mathbin{\upharpoonright}L=y\mathbin{\upharpoonright}L
  \quad\Longrightarrow\quad
  \Phi(x)\mathbin{\upharpoonright}N
   =\Phi(y)\mathbin{\upharpoonright}N.
\]
If $\ell\geq L$, all vertices of a canonical level-$\ell$ source
cycle agree below $L$.  Their images therefore agree below $N$.
Every target edge in the resulting closed walk has level at least
$N$.  Lemma~\ref{lem:rxcb-level-counting} now gives the conclusion.
\end{proof}

Even Lemma~\ref{lem:rxcb-tail-monoid} is not enough for a general
diagonalization.  For example, if $f(\ell)=\prod_{i=2}^{\ell+2}i$, then for every
$g$ and every fixed $N$, all sufficiently large $f(\ell)$ are
multiples of $g(N)$ and hence belong to the monoid generated by the
$N$th tail of $g$.  A successful proof for ordinary Borel
homomorphisms must therefore use the incidence pattern of canonical
cycles, rather than only global or local cycle-length semigroups.

\subsection{A direct ordinary-homomorphism lower bound}
\label{subsec:rxcb-bounded}

\begin{proposition}[Bounded sources]
\label{prop:rxcb-bounded-source}
Suppose that $f$ is bounded, $E$ is dense, and $g$ is one-to-one.
Then there is no Borel homomorphism from $\mathcal G_f$ to
$\mathcal G_g$.
\end{proposition}

\begin{proof}
Let $K$ bound $f$, and suppose that
$\Phi:\mathcal G_f\to\mathcal G_g$ is a Borel homomorphism.  Every
source edge belongs to a canonical cycle of at most $K$ edges.  Its
image therefore lies on a directed closed walk of at most $K$ edges.
By Lemma~\ref{lem:rxcb-level-counting}, every image edge has level in
\[
  M=\{m:g(m)\leq K\}.
\]
The set $M$ is finite because $g$ is one-to-one.  The map
\[
  c(x)=\langle x(m):m\in M\rangle
\]
takes finitely many values and properly colors the subdigraph of
$\mathcal G_g$ formed by the edges whose levels lie in $M$.
Consequently, $c\circ\Phi$ is a finite Borel proper coloring of
$\mathcal G_f$, and hence a finite Borel dicoloring.  This contradicts
Lemma~3.8 of \cite{RX24}.
\end{proof}

\begin{corollary}
\label{cor:rxcb-no-uniform-basis}
No countable family of $k$-uniform Borel digraphs, where $k$ may vary
from one member to another, is a basis for $\mathcal C_\infty$ under
ordinary Borel homomorphism.
\end{corollary}

\begin{proof}
Suppose that $\{H_n:n\in\omega\}$ were such a basis.  Apply the
boundedness clause of Theorem~3.25 in \cite{RX24} to choose, for every
$n$, a bounded $f_n$ and a continuous homomorphism
\[
  \mathcal G_{f_n}\longrightarrow H_n.
\]
Let $g$ be one-to-one.  By Lemma~3.8,
$\mathcal G_g\in\mathcal C_\infty$.  Hence some $H_n$ admits a Borel
homomorphism to $\mathcal G_g$.  Composition gives a Borel
homomorphism $\mathcal G_{f_n}\to\mathcal G_g$, contrary to
Proposition~\ref{prop:rxcb-bounded-source}.
\end{proof}

\subsection{A complete diagonalization for minimal homomorphisms}
\label{subsec:rxcb-minimal}

\begin{lemma}[Rigidity of simple cycles]
\label{lem:rxcb-cycle-rigidity}
Every directed cycle with no repeated vertices in $\mathcal G_g$ is a
canonical cycle at one level $m$.  In particular, every minimal cycle
in $\mathcal G_g$ has exactly $g(m)$ vertices for some $m$.
\end{lemma}

\begin{proof}
Let $C$ be a directed cycle with no repeated vertices, and let $m$ be
the largest level of an edge of $C$.  Every edge of $C$ preserves all
coordinates above $m$.  Thus all level-$m$ edges of $C$ have the same
tail above $m$ and belong to one canonical level-$m$ cycle.  By
Lemma~\ref{lem:rxcb-level-counting}, the number of level-$m$ edges in
$C$ is a positive multiple of $g(m)$.  Since $C$ has no repeated
vertices, it can use at most the $g(m)$ distinct level-$m$ edges with
that tail.  It therefore uses all of them exactly once.

The terminal vertex of each level-$m$ edge is the initial vertex of
the next level-$m$ edge in the canonical cyclic order.  A path made of
lower-level edges preserves coordinate $m$ and the entire tail above
$m$.  Hence a nonempty lower-level segment between two consecutive
level-$m$ edges would begin and end at the same vertex, contradicting
simplicity of $C$.  No lower-level edge occurs, so $C$ is the canonical
level-$m$ cycle.
\end{proof}

\begin{lemma}[Diagonal function]
\label{lem:rxcb-minimal-diagonal}
For every sequence $\langle f_n:n\in\omega\rangle$ in
$(\omega\setminus2)^\omega$, there is a one-to-one
$g\in(\omega\setminus2)^\omega$ such that
\[
  \mathcal G_{f_n}\not\preceq^{\min}_{\mathrm B}\mathcal G_g
  \quad\text{for every }n.
\]
\end{lemma}

\begin{proof}
Let
\[
  I=\{n:\operatorname{ran}(f_n)\text{ is infinite}\}.
\]
Enumerate $I$ without repetition as $\langle n_j:j\in J\rangle$,
where $J$ is either a finite initial segment of $\omega$ or all of
$\omega$.  Recursively choose numbers $a_j,b_j\geq2$, for $j\in J$,
so that all chosen numbers are distinct and
\[
  a_j\in\operatorname{ran}(f_{n_j}).
\]
At stage $j$, first choose a number $b_j$ not used earlier.  Then use
the infinitude of $\operatorname{ran}(f_{n_j})$ to choose $a_j$
outside all numbers chosen so far, including $b_j$.  If $J$ is finite,
add infinitely many new numbers to $\{b_j:j\in J\}$.  In either case,
obtain an infinite set $B\subseteq\omega\setminus2$ such that
\[
  a_j\notin B\quad\text{for every }j\in J.
\]
Let $g$ be a one-to-one enumeration of $B$.

Suppose first that $n=n_j\in I$.  A minimal homomorphism would map a
canonical $a_j$-cycle in $\mathcal G_{f_n}$ to a minimal cycle with the
same number of vertices.  Lemma~\ref{lem:rxcb-cycle-rigidity} would
then give $a_j=g(m)$ for some $m$, contrary to $a_j\notin B$.

If $n\notin I$, then $f_n$ has finite range and is bounded.
Proposition~\ref{prop:rxcb-bounded-source} rules out even an ordinary
Borel homomorphism from $\mathcal G_{f_n}$ to $\mathcal G_g$.
\end{proof}

\begin{theorem}
\label{thm:rxcb-no-countable-minimal-basis}
There is no countable basis for $\mathcal C_\infty$ under Borel
minimal homomorphism.  The same conclusion holds when all
homomorphisms in the definition of basis are required to be continuous
and minimal.
\end{theorem}

\begin{proof}
Suppose that $\{H_n:n\in\omega\}$ is a basis under Borel minimal
homomorphism.  By Theorem~3.23 of \cite{RX24}, choose, for every $n$, a
function $f_n$ and a continuous minimal homomorphism
\[
  \alpha_n:\mathcal G_{f_n}\longrightarrow H_n.
\]
Apply Lemma~\ref{lem:rxcb-minimal-diagonal} and obtain a one-to-one $g$
such that no $\mathcal G_{f_n}$ admits a Borel minimal homomorphism to
$\mathcal G_g$.  Lemma~3.8 gives
$\mathcal G_g\in\mathcal C_\infty$.  Hence the basis assumption yields
an $n$ and a Borel minimal homomorphism
\[
  \beta:H_n\longrightarrow\mathcal G_g.
\]
Minimal homomorphisms are closed under composition: the first map
sends each induced directed cycle isomorphically onto an induced cycle
of the same size, to which the second map's minimality applies.
Therefore $\beta\circ\alpha_n$ is a Borel minimal homomorphism from
$\mathcal G_{f_n}$ to $\mathcal G_g$, a contradiction.  The continuous
version is identical.
\end{proof}

\subsection{Status of the ordinary-homomorphism conjecture}
\label{subsec:rxcb-status}

Theorem~\ref{thm:rxcb-no-countable-minimal-basis} does not settle the
basis problem under ordinary Borel homomorphism.  This is not merely a
regularity gap.  An ordinary homomorphism may fold a canonical cycle
or send it around several target cycles, and
Lemma~\ref{lem:rxcb-level-counting} shows that the first numerical
invariant is the additive monoid generated by target cycle sizes.
The family
\[
  \bigl\{\{k,2k,3k,\dots\}:k\geq2\bigr\}
\]
shows that no range-only diagonal can handle every countable family,
and factorial sequences defeat even the tail-monoid obstruction for
continuous maps.

Accordingly, a proof of the full conjecture must exploit the way
canonical cycles at different levels intersect, or prove a
canonization theorem saying that every Borel homomorphism into a
suitably chosen $\mathcal G_0(E,g)$ becomes cycle-injective or minimal
on a canonical substructure.  No such statement is proved in
\cite{RX24}, and the arguments above do not establish it.  On the
other hand, Corollary~\ref{cor:rxcb-conditional-full} shows that a
verified pairwise-incomparability theorem for an actual cofinal
Raghavan--Xiao family would immediately prove the ordinary-homomorphism
conjecture.
